\begin{lemma}[Restriction of a weak gradient to an open subset]
Let $n\ge1$ and $q\ge1$, let $\Omega,\Omega'\subseteq\R^n$ be open with
$\Omega'\subseteq\Omega$, let $f\colon\R^n\to\R$ and let
$g\colon\R^n\to\R^n$.  Then:
\begin{enumerate}
\item if $f\in W^{1,q}(\Omega)$ with weak gradient $g$ (\noderef{MemW1pOn}),
  then $f\in W^{1,q}(\Omega')$ with weak gradient $g$;
\item if $f\in W^{1,q}_{\mathrm{loc}}(\Omega)$ with weak gradient $g$
  (\noderef{MemW1pLocOn}), then $f\in W^{1,q}_{\mathrm{loc}}(\Omega')$ with weak
  gradient $g$; and
\item if $f\in W^{1,q}(\Omega)$ with weak gradient $g$, then
  $f\in W^{1,q}_{\mathrm{loc}}(\Omega)$ with weak gradient $g$; and
\item if $f\in W^{1,q}_{\mathrm{loc}}(\Omega)$ with weak gradient $g$ and
  $\overline{\Omega'}$ is compact with $\overline{\Omega'}\subseteq\Omega$, then
  $f\in W^{1,q}(\Omega')$ with weak gradient $g$.
\end{enumerate}

The four clauses are the four ways the development passes between the two
predicates and the two sets.  Clause (3) is not a special case of (1) or (2):
it changes the predicate, not the set, and it is what lets a global
$W^{1,q}$ hypothesis feed a node whose statement asks for the local one.
Clause (4) is the converse passage, and it is the one that costs something: a
local hypothesis on $\Omega$ yields a \emph{global} $W^{1,q}$ conclusion on
$\Omega'$ only when $\Omega'$ is compactly contained in $\Omega$, which is
exactly the hypothesis that makes $\overline{\Omega'}$ an admissible compact set
for the local clause.  It is the standard exhaustion step --- a solution known
only locally on a ball is in $W^{1,q}$ of every smaller concentric ball --- and
without it as a statement the step gets written out inline at each use.  In
clauses (1) and (2) the weak gradient is unchanged --- it is the \emph{same}
field $g$, not a restriction of it --- because the defining identity of
\noderef{HasWeakGradientOn} on a smaller set tests $g$ against fewer functions
and asks nothing of $g$ off $\Omega'$.  The converse of (1) is false in general
and is not claimed: a function may be in $W^{1,q}$ of each of two open sets
without being in $W^{1,q}$ of their union; the opposite direction, extending a
compactly supported Sobolev function from a smaller set to a larger one, is a
separate statement with its own hypotheses and is not part of this node.

As with the other facts in this family the statement is abstract: no operator,
no solution, no structure condition.
\end{lemma}
\begin{proof}
Throughout, $D_i$ is the derivative in the direction of the $i$-th standard
basis vector, and a test function for an open set $U$ means a function
$\varphi\in C^\infty(\R^n)$ with compact support and
$\operatorname{supp}\varphi\subseteq U$, the test class of
\noderef{HasWeakGradientOn}.

\emph{Step 1: the weak-gradient identity passes to $\Omega'$.}  Assume
$\mathrm{HasWeakGradientOn}(n,\Omega,f,g)$ (\noderef{HasWeakGradientOn}), let
$\varphi$ be a test function for $\Omega'$ and let $i\in\{1,\dots,n\}$.  Since
$\operatorname{supp}\varphi\subseteq\Omega'\subseteq\Omega$, the function
$\varphi$ is also a test function for $\Omega$, so the hypothesis gives
\[
\int_\Omega f\,D_i\varphi\,dx=-\int_\Omega g^i\,\varphi\,dx .
\tag{1}
\]
Both integrands vanish off $\operatorname{supp}\varphi$: the function $\varphi$
vanishes on the open complement of its support, hence so does $D_i\varphi$.
Therefore, as functions on $\R^n$,
\[
\mathbf1_\Omega\cdot\bigl(f\,D_i\varphi\bigr)
=\mathbf1_{\Omega'}\cdot\bigl(f\,D_i\varphi\bigr),
\qquad
\mathbf1_\Omega\cdot\bigl(g^i\varphi\bigr)
=\mathbf1_{\Omega'}\cdot\bigl(g^i\varphi\bigr),
\]
because the two indicators agree on $\operatorname{supp}\varphi$ --- both are
$1$ there, as $\operatorname{supp}\varphi\subseteq\Omega'\subseteq\Omega$ ---
and both products vanish off $\operatorname{supp}\varphi$.  An integral over a
set depends on the integrand only through its product with the indicator of that
set, so each side of (1) is unchanged when $\Omega$ is replaced by $\Omega'$:
\[
\int_{\Omega'}f\,D_i\varphi\,dx=-\int_{\Omega'}g^i\,\varphi\,dx .
\]
As $\varphi$ and $i$ were arbitrary, this is
$\mathrm{HasWeakGradientOn}(n,\Omega',f,g)$.  No integrability was used: the two
displayed equalities are equalities of integrands.

\emph{Step 2: clause (1).}  Assume $f\in W^{1,q}(\Omega)$ with weak gradient $g$.
Its first clause gives $\mathrm{HasWeakGradientOn}(n,\Omega,f,g)$, so Step 1
gives $\mathrm{HasWeakGradientOn}(n,\Omega',f,g)$.  Its second and third clauses
give $f\in L^q(\Omega)$ and $|g|\in L^q(\Omega)$, and since
$\Omega'\subseteq\Omega$ the Preamble's transport fact for $L^q$ membership
under restriction to a subset gives $f\in L^q(\Omega')$ and
$|g|\in L^q(\Omega')$.  By \noderef{MemW1pOn}, $f\in W^{1,q}(\Omega')$ with weak
gradient $g$.

\emph{Step 3: clause (2).}  Assume $f\in W^{1,q}_{\mathrm{loc}}(\Omega)$ with
weak gradient $g$.  Its first clause and Step 1 give
$\mathrm{HasWeakGradientOn}(n,\Omega',f,g)$.  For the local clause on $\Omega'$,
let $K\subseteq\Omega'$ be compact; then $K\subseteq\Omega$, so the local clause
of \noderef{MemW1pLocOn} on $\Omega$ applies to $K$ verbatim and gives
$f\in L^q(K)$ and $|g|\in L^q(K)$.  By \noderef{MemW1pLocOn},
$f\in W^{1,q}_{\mathrm{loc}}(\Omega')$ with weak gradient $g$.  (Here nothing
had to be transported: the local clause quantifies over compact subsets, and a
compact subset of $\Omega'$ is one of $\Omega$.)

\emph{Step 4': clause (4).}  Assume $f\in W^{1,q}_{\mathrm{loc}}(\Omega)$ with
weak gradient $g$, and let $\overline{\Omega'}$ be compact with
$\overline{\Omega'}\subseteq\Omega$.  The first clause and Step 1 give
$\mathrm{HasWeakGradientOn}(n,\Omega',f,g)$ (the hypothesis
$\Omega'\subseteq\Omega$ of the lemma is what Step 1 uses, and it also follows
from $\overline{\Omega'}\subseteq\Omega$).  For the $L^q$ clauses, apply the
local clause of \noderef{MemW1pLocOn} to the compact set
$\overline{\Omega'}\subseteq\Omega$: it gives $f\in L^q(\overline{\Omega'})$ and
$|g|\in L^q(\overline{\Omega'})$.  Since $\Omega'\subseteq\overline{\Omega'}$,
the Preamble's transport fact for $L^q$ membership under restriction to a subset
gives $f\in L^q(\Omega')$ and $|g|\in L^q(\Omega')$.  By \noderef{MemW1pOn},
$f\in W^{1,q}(\Omega')$ with weak gradient $g$.

\emph{Step 4: clause (3).}  Assume $f\in W^{1,q}(\Omega)$ with weak gradient
$g$.  Its first clause is already
$\mathrm{HasWeakGradientOn}(n,\Omega,f,g)$, which is the first clause of
\noderef{MemW1pLocOn} on the same set.  For the local clause, let
$K\subseteq\Omega$ be compact; from $f\in L^q(\Omega)$ and $|g|\in L^q(\Omega)$
the Preamble's transport fact for restriction to a subset gives $f\in L^q(K)$
and $|g|\in L^q(K)$.  By \noderef{MemW1pLocOn},
$f\in W^{1,q}_{\mathrm{loc}}(\Omega)$ with weak gradient $g$.
\end{proof}
